\documentclass[11pt,a4paper]{article}
\usepackage[utf8]{inputenc}
\usepackage[T1]{fontenc}
\usepackage[english]{babel}
\usepackage{amsmath,amssymb,amsthm}
\usepackage{booktabs}
\usepackage{enumitem}
\usepackage[margin=2.5cm]{geometry}
\usepackage{xcolor}
\usepackage{hyperref}

\newcommand{\vect}[1]{\mathbf{#1}}
\newcommand{\norm}[1]{\lVert #1 \rVert}
\newcommand{\code}[1]{\texttt{#1}}
\DeclareMathOperator{\atantwo}{atan2}
\DeclareMathOperator{\clamp}{clamp}

\newenvironment{idea}
  {\par\medskip\noindent\begingroup\color{blue!50!black}\textbf{Key idea.}\ }
  {\endgroup\par\medskip}

\title{Diamond Tracking Mathematics\\\large (4 $\to$ 3 $\to$ 2 LEDs)}
\author{\code{wiimote-pico-mouse}}
\date{}

\begin{document}
\maketitle

\begin{abstract}
This document explains the geometric pipeline implemented in
\code{src/diamond\_tracker.c} and \code{src/one\_euro\_filter.c}. Its goal is
to convert the infrared detections from the Wiimote's PixArt camera into a
jitter-free absolute USB HID pointer, without cursor jumps when LEDs are lost.
\end{abstract}

\tableofcontents

\section{Overview}

The standard sensor bar is replaced by four IR LEDs placed at the vertices of a
diamond around the TV: \code{Top}, \code{Right}, \code{Bottom}, and
\code{Left}. The Wiimote's camera sees these four points; if we know
\emph{where} they land in the camera image, we can deduce \emph{where} on the
screen the controller is aiming.

In each frame, the pipeline runs the following steps:
\begin{enumerate}[nosep]
  \item Center the camera coordinates and measure the controller's roll angle
        from the accelerometer (Section~\ref{sec:camera}).
  \item Depending on how many LEDs are visible, reconstruct the four diamond
        vertices: 4 LEDs (Section~\ref{sec:four}), 3 LEDs
        (Section~\ref{sec:three}), or 2 LEDs (Section~\ref{sec:two}).
  \item Compute the screen point corresponding to the camera's optical center
        using a 2D projective homography.
  \item Apply the 2-point affine calibration (Section~\ref{sec:calib}) and
        smooth the output with the adaptive 1-Euro filter
        (Section~\ref{sec:filter}).
\end{enumerate}

%------------------------------------------------------------------
\section{Camera Coordinate Frame and Roll}
\label{sec:camera}

\subsection{Coordinate Centering}

The PixArt IR camera outputs up to 4 valid blobs on a $1024 \times 768$ pixel
grid whose optical center lies at $(512, 384)$. To align the axes with the
screen ($+u$ right, $+v$ down), each raw detection $(x_{\text{raw}},
y_{\text{raw}})$ is centered as:
\begin{equation}
  u = 512 - x_{\text{raw}}, \qquad v = y_{\text{raw}} - 384 .
\end{equation}
With this convention, the optical axis (where the Wiimote is aimed) always sits
at the origin $(u,v) = (0,0)$.

\begin{idea}
The entire tracking problem reduces to one question: \emph{where inside the
detected diamond does the origin $(0,0)$ fall?} That relative position is the
point on the screen we are aiming at.
\end{idea}

\subsection{Controller Roll}

The 10-bit ADXL330 accelerometer (zero-centered at \code{512} in report
\code{0x33}) measures the gravity vector $(a_x, a_y, a_z)$, from which the
controller's roll around its aiming axis is computed:
\begin{equation}
  \theta_{\text{roll}} = \atantwo(a_x, a_z) .
\end{equation}
To rotate a camera vector $\vect{p} = (u,v)$ back into the upright TV plane,
we apply the 2D rotation matrix:
\begin{equation}
  R(\theta_{\text{roll}}) =
  \begin{bmatrix}
    \cos\theta_{\text{roll}} & -\sin\theta_{\text{roll}} \\
    \sin\theta_{\text{roll}} &  \cos\theta_{\text{roll}}
  \end{bmatrix}.
\end{equation}

%------------------------------------------------------------------
\section{4-LED State: Full 2D Projective Homography}
\label{sec:four}

When all four diamond vertices are visible, there is enough information to
solve the perspective mapping exactly.

\subsection{Vertex Identification}

The camera reports unlabeled blobs, so we must determine which blob is
\code{Top}, \code{Right}, \code{Bottom}, and \code{Left}. All $4! = 24$
permutations are evaluated, rejecting self-intersecting quadrilaterals by
requiring:
\begin{equation}
  (P_R - P_L) \times (P_B - P_T) > 0 ,
\end{equation}
where $\times$ denotes the 2D cross product (determinant). Each candidate
permutation is scored using two criteria:
\begin{itemize}[nosep]
  \item \textbf{Orientation alignment:} centroid-relative vectors, derotated by
        $R(\theta_{\text{roll}})$, should match the canonical directions
        $\{(0,-1),\,(1,0),\,(0,1),\,(-1,0)\}$.
  \item \textbf{Temporal continuity:} vertices should stay close to the
        previous frame's vertices rotated by $-\Delta\theta_{\text{roll}}$ to
        compensate for hand rotation.
\end{itemize}

\subsection{Homography Computation}

A homography is a projective transformation of the plane, represented by a
$3\times 3$ matrix $H$ (with $h_8 = 1$ to fix scale) that maps the four
camera vertices $P_k = (u_k, v_k)$ to the normalized screen edge midpoints
$(X_k, Y_k)$:

\begin{center}
\begin{tabular}{lc}
\toprule
Vertex & Normalized Screen Target $(X_k, Y_k)$ \\
\midrule
\code{Top}    ($k=0$) & $(0.5,\ 0.0)$ \\
\code{Right}  ($k=1$) & $(1.0,\ 0.5)$ \\
\code{Bottom} ($k=2$) & $(0.5,\ 1.0)$ \\
\code{Left}   ($k=3$) & $(0.0,\ 0.5)$ \\
\bottomrule
\end{tabular}
\end{center}

For any camera point $(u,v)$, the mapping is:
\begin{equation}
  X = \frac{h_0 u + h_1 v + h_2}{h_6 u + h_7 v + 1}, \qquad
  Y = \frac{h_3 u + h_4 v + h_5}{h_6 u + h_7 v + 1}.
\end{equation}
Clearing the denominator, each correspondence $P_k \mapsto (X_k,Y_k)$ yields
\emph{two linear equations} in the unknowns $h_0,\dots,h_7$:
\begin{align}
  u_k h_0 + v_k h_1 + h_2 - X_k u_k h_6 - X_k v_k h_7 &= X_k, \\
  u_k h_3 + v_k h_4 + h_5 - Y_k u_k h_6 - Y_k v_k h_7 &= Y_k.
\end{align}
The 4 vertices form an $8\times 8$ linear system, which is solved in double
precision using Gaussian elimination with partial pivoting (after coordinate
normalization for numerical stability).

\subsection{Aiming Point}

Because the optical axis sits at $(u,v) = (0,0)$, the denominator evaluates to
$1$ and all terms involving $u$ or $v$ vanish, giving directly:
\begin{equation}
  X_{\text{raw}} = h_2, \qquad Y_{\text{raw}} = h_5 .
\end{equation}

\subsection{Live Geometry State Update}

Every 4-LED frame updates:
\begin{itemize}[nosep]
  \item The diagonals $\vect{d}_H = P_R - P_L$ and $\vect{d}_V = P_B - P_T$,
        their lengths $L_H = \norm{\vect{d}_H}$, $L_V = \norm{\vect{d}_V}$, and
        the aspect ratio $\rho = L_H / L_V$.
  \item The diagonal midpoints $M_H = (P_L + P_R)/2$ and
        $M_V = (P_T + P_B)/2$, and the perspective midpoint offset
        \begin{equation}
          \boldsymbol{\delta} = M_V - M_H
            = \alpha_H \vect{d}_H + \alpha_V \vect{d}_V ,
        \end{equation}
        which captures how much the two diagonal centers separate due to
        perspective (for a perfectly fronto-parallel diamond,
        $\boldsymbol\delta = 0$).
  \item A drift-free reference template \code{ref\_vertices[4]} and its roll
        angle \code{ref\_roll}.
\end{itemize}

%------------------------------------------------------------------
\section{3-LED State: Perspective-Corrected 4th Vertex Reconstruction}
\label{sec:three}

When one LED exits the camera field of view or is occluded, 3 blobs remain. The
idea is to \emph{reconstruct} the missing vertex from the learned geometry and
feed all 4 vertices into the exact same homography, preventing cursor jumps.

All $4 \times 3! = 24$ assignments (\code{missing\_slot}, 3-point permutation)
are evaluated.

\paragraph{Case A: \code{Top} or \code{Bottom} is missing.}
The horizontal diagonal $\vect{d}_H = P_R - P_L$ is complete. The vertical
diagonal is approximated by rotating $\vect{d}_H$ by $90^\circ$ and scaling by
the aspect ratio:
\begin{equation}
  \vect{d}_{V,\text{approx}} = \frac{(-d_{H,v},\ d_{H,u})}{\rho}.
\end{equation}
The perspective-corrected vertical midpoint is:
\begin{equation}
  M_V = M_H + \alpha_H \vect{d}_H + \alpha_V \vect{d}_{V,\text{approx}},
\end{equation}
and the missing vertex is reconstructed by reflecting its opposite vertex across
$M_V$:
\begin{equation}
  P_{\text{missing}} = 2 M_V - P_{\text{opposite}} .
\end{equation}

\paragraph{Case B: \code{Left} or \code{Right} is missing.}
Symmetrically, from the complete vertical diagonal $\vect{d}_V = P_B - P_T$:
\begin{equation}
  \vect{d}_{H,\text{approx}} = (d_{V,v},\ -d_{V,u})\cdot\rho ,
\end{equation}
\begin{equation}
  M_H = M_V - \alpha_H \vect{d}_{H,\text{approx}} - \alpha_V \vect{d}_V,
  \qquad
  P_{\text{missing}} = 2 M_H - P_{\text{opposite}} .
\end{equation}

\paragraph{Selecting the best candidate.}
The winning hypothesis minimizes a combination of:
\begin{itemize}[nosep]
  \item translation-invariant pairwise edge error against the roll-updated
        previous frame;
  \item spatial continuity of the 3 visible dots;
  \item a preference for reconstructions that fall \emph{outside} the sensor
        rectangle (as expected when an LED leaves the camera FOV).
\end{itemize}

%------------------------------------------------------------------
\section{2-LED State: 2D Similarity Fallback with Accelerometer Roll}
\label{sec:two}

On large TVs or at short distances, aiming into a corner can leave only 2
adjacent LEDs visible (e.g., \code{Top} and \code{Right}). Two points cannot
define a projective homography, but they uniquely define a \emph{2D similarity
transform} (uniform scale + rotation + translation).

\begin{enumerate}
  \item \textbf{Rotated template.} The stored \code{ref\_vertices[4]} template
        (captured during the last 4-LED frame) is rotated to the live
        accelerometer roll via $\Delta\theta = \theta_{\text{roll}} -
        \theta_{\text{ref}}$, yielding vectors $\vect{r}_k$.
  \item \textbf{Pair identification.} All $4\times 3 = 12$ ordered pairs
        $(a,b)$ are scored by comparing the observed vector
        $\vect{o} = P_1 - P_0$ against the template vector
        $\vect{e} = \vect{r}_b - \vect{r}_a$ in angle, scale, and temporal
        proximity to \code{last\_vertices}.
  \item \textbf{Exact similarity transform.} For the winning pair $(a,b)$, we
        find the similarity mapping $\vect{r}_a \mapsto P_0$ and
        $\vect{r}_b \mapsto P_1$. Written as a complex number $c + i d$
        (encoding both scale and rotation), its coefficients are:
        \begin{equation}
          c = \frac{\vect{o}\cdot\vect{e}}{\norm{\vect{e}}^2}, \qquad
          d = \frac{o_v e_u - o_u e_v}{\norm{\vect{e}}^2}.
        \end{equation}
        All 4 vertices are then reconstructed as:
        \begin{equation}
          V_k = \frac{P_0 + P_1}{2} +
          \begin{bmatrix} c & -d \\ d & c \end{bmatrix}
          \left(\vect{r}_k - \frac{\vect{r}_a + \vect{r}_b}{2}\right).
        \end{equation}
  \item \textbf{Zero drift.} Because \code{ref\_vertices} is \emph{never}
        overwritten in 2-LED mode, no geometric drift accumulates even when
        holding the pointer in a corner indefinitely.
\end{enumerate}

The 4 reconstructed vertices $V_k$ are then passed through the same homography
described in Section~\ref{sec:four}.

%------------------------------------------------------------------
\section{1-LED and 0-LED States}

When fewer than 2 IR LEDs are visible, there is not enough geometric
information to track the pointer. The cursor freezes smoothly at its last valid
coordinates while all Wiimote buttons remain operational.

%------------------------------------------------------------------
\section{2-Point Affine Calibration (Flash)}
\label{sec:calib}

To compensate for bezel geometry and mounting tolerances, calibration mode
(hold \code{+} and \code{-} for 3\,s) locks the USB cursor to two known screen
targets:
\begin{equation}
  T_1 = (0.10,\ 0.10), \qquad T_2 = (0.90,\ 0.90),
\end{equation}
and records the raw homography outputs $P_1$ and $P_2$ when \code{A} is
pressed at each target. For the $X$ axis, a linear mapping
$X_{\text{cal}} = s_x X_{\text{raw}} + b_x$ passing through both points is
computed as:
\begin{equation}
  s_x = \frac{0.90 - 0.10}{P_{2,x} - P_{1,x}}, \qquad
  b_x = 0.10 - s_x P_{1,x},
\end{equation}
\begin{equation}
  X_{\text{cal}} = \clamp_{[0,1]}\!\left(s_x X_{\text{raw}} + b_x\right).
\end{equation}
The $Y$ axis is calibrated analogously. All coefficients and reference
diagonals are stored in Flash, protected by a magic header and an IEEE~802.3
CRC-32.

%------------------------------------------------------------------
\section{Adaptive 1-Euro Filter}
\label{sec:filter}

Natural hand tremor becomes noticeable at distances $\ge 2.5$\,m. A fixed
low-pass filter would remove the tremor but introduce lag during fast
movements. The 1-Euro filter solves this trade-off by adapting its cutoff
frequency to the pointer's speed: \emph{low speed $\Rightarrow$ heavy
smoothing; high speed $\Rightarrow$ minimal lag}.
It is applied per axis to $(X_{\text{cal}}, Y_{\text{cal}})$
(\code{src/one\_euro\_filter.c}). Let $x_i$ be the current sample and $\Delta t$
the time between samples.

\paragraph{Smoothing factor.} A first-order low-pass filter with cutoff
frequency $f$ uses:
\begin{equation}
  \alpha(\Delta t, f) = \frac{2\pi f \Delta t}{1 + 2\pi f \Delta t},
  \qquad
  \mathrm{LPF}(x, \alpha) = \alpha\, x + (1-\alpha)\, x_{\text{prev}} .
\end{equation}

\paragraph{Step 1: smoothed velocity.} The velocity is estimated and filtered
with a fixed derivative cutoff $f_d = 1.0$\,Hz:
\begin{equation}
  v_{\text{filt},i} = \mathrm{LPF}\!\left(
    \frac{x_i - x_{\text{filt},i-1}}{\Delta t},\ \alpha(\Delta t, f_d)\right).
\end{equation}

\paragraph{Step 2: adaptive cutoff.} The cutoff frequency increases linearly
with speed ($f_{c,\min} = 1.2$\,Hz and $\beta = 2.5$ for normalized $[0,1]$
coordinates):
\begin{equation}
  f_c = f_{c,\min} + \beta\,\lvert v_{\text{filt},i}\rvert .
\end{equation}

\paragraph{Step 3: filtered position.}
\begin{equation}
  x_{\text{filt},i} = \mathrm{LPF}\!\left(x_i,\ \alpha(\Delta t, f_c)\right).
\end{equation}
This $x_{\text{filt},i}$ is the final value sent to the absolute USB HID
pointer.

\end{document}
